% Options for packages loaded elsewhere
\PassOptionsToPackage{unicode}{hyperref}
\PassOptionsToPackage{hyphens}{url}
\documentclass[
]{article}
\usepackage{xcolor}
\usepackage[margin=1in]{geometry}
\usepackage{amsmath,amssymb}
\setcounter{secnumdepth}{-\maxdimen} % remove section numbering
\usepackage{iftex}
\ifPDFTeX
  \usepackage[T1]{fontenc}
  \usepackage[utf8]{inputenc}
  \usepackage{textcomp} % provide euro and other symbols
\else % if luatex or xetex
  \usepackage{unicode-math} % this also loads fontspec
  \defaultfontfeatures{Scale=MatchLowercase}
  \defaultfontfeatures[\rmfamily]{Ligatures=TeX,Scale=1}
\fi
\usepackage{lmodern}
\ifPDFTeX\else
  % xetex/luatex font selection
\fi
% Use upquote if available, for straight quotes in verbatim environments
\IfFileExists{upquote.sty}{\usepackage{upquote}}{}
\IfFileExists{microtype.sty}{% use microtype if available
  \usepackage[]{microtype}
  \UseMicrotypeSet[protrusion]{basicmath} % disable protrusion for tt fonts
}{}
\makeatletter
\@ifundefined{KOMAClassName}{% if non-KOMA class
  \IfFileExists{parskip.sty}{%
    \usepackage{parskip}
  }{% else
    \setlength{\parindent}{0pt}
    \setlength{\parskip}{6pt plus 2pt minus 1pt}}
}{% if KOMA class
  \KOMAoptions{parskip=half}}
\makeatother
\usepackage{color}
\usepackage{fancyvrb}
\newcommand{\VerbBar}{|}
\newcommand{\VERB}{\Verb[commandchars=\\\{\}]}
\DefineVerbatimEnvironment{Highlighting}{Verbatim}{commandchars=\\\{\}}
% Add ',fontsize=\small' for more characters per line
\usepackage{framed}
\definecolor{shadecolor}{RGB}{248,248,248}
\newenvironment{Shaded}{\begin{snugshade}}{\end{snugshade}}
\newcommand{\AlertTok}[1]{\textcolor[rgb]{0.94,0.16,0.16}{#1}}
\newcommand{\AnnotationTok}[1]{\textcolor[rgb]{0.56,0.35,0.01}{\textbf{\textit{#1}}}}
\newcommand{\AttributeTok}[1]{\textcolor[rgb]{0.13,0.29,0.53}{#1}}
\newcommand{\BaseNTok}[1]{\textcolor[rgb]{0.00,0.00,0.81}{#1}}
\newcommand{\BuiltInTok}[1]{#1}
\newcommand{\CharTok}[1]{\textcolor[rgb]{0.31,0.60,0.02}{#1}}
\newcommand{\CommentTok}[1]{\textcolor[rgb]{0.56,0.35,0.01}{\textit{#1}}}
\newcommand{\CommentVarTok}[1]{\textcolor[rgb]{0.56,0.35,0.01}{\textbf{\textit{#1}}}}
\newcommand{\ConstantTok}[1]{\textcolor[rgb]{0.56,0.35,0.01}{#1}}
\newcommand{\ControlFlowTok}[1]{\textcolor[rgb]{0.13,0.29,0.53}{\textbf{#1}}}
\newcommand{\DataTypeTok}[1]{\textcolor[rgb]{0.13,0.29,0.53}{#1}}
\newcommand{\DecValTok}[1]{\textcolor[rgb]{0.00,0.00,0.81}{#1}}
\newcommand{\DocumentationTok}[1]{\textcolor[rgb]{0.56,0.35,0.01}{\textbf{\textit{#1}}}}
\newcommand{\ErrorTok}[1]{\textcolor[rgb]{0.64,0.00,0.00}{\textbf{#1}}}
\newcommand{\ExtensionTok}[1]{#1}
\newcommand{\FloatTok}[1]{\textcolor[rgb]{0.00,0.00,0.81}{#1}}
\newcommand{\FunctionTok}[1]{\textcolor[rgb]{0.13,0.29,0.53}{\textbf{#1}}}
\newcommand{\ImportTok}[1]{#1}
\newcommand{\InformationTok}[1]{\textcolor[rgb]{0.56,0.35,0.01}{\textbf{\textit{#1}}}}
\newcommand{\KeywordTok}[1]{\textcolor[rgb]{0.13,0.29,0.53}{\textbf{#1}}}
\newcommand{\NormalTok}[1]{#1}
\newcommand{\OperatorTok}[1]{\textcolor[rgb]{0.81,0.36,0.00}{\textbf{#1}}}
\newcommand{\OtherTok}[1]{\textcolor[rgb]{0.56,0.35,0.01}{#1}}
\newcommand{\PreprocessorTok}[1]{\textcolor[rgb]{0.56,0.35,0.01}{\textit{#1}}}
\newcommand{\RegionMarkerTok}[1]{#1}
\newcommand{\SpecialCharTok}[1]{\textcolor[rgb]{0.81,0.36,0.00}{\textbf{#1}}}
\newcommand{\SpecialStringTok}[1]{\textcolor[rgb]{0.31,0.60,0.02}{#1}}
\newcommand{\StringTok}[1]{\textcolor[rgb]{0.31,0.60,0.02}{#1}}
\newcommand{\VariableTok}[1]{\textcolor[rgb]{0.00,0.00,0.00}{#1}}
\newcommand{\VerbatimStringTok}[1]{\textcolor[rgb]{0.31,0.60,0.02}{#1}}
\newcommand{\WarningTok}[1]{\textcolor[rgb]{0.56,0.35,0.01}{\textbf{\textit{#1}}}}
\usepackage{graphicx}
\makeatletter
\newsavebox\pandoc@box
\newcommand*\pandocbounded[1]{% scales image to fit in text height/width
  \sbox\pandoc@box{#1}%
  \Gscale@div\@tempa{\textheight}{\dimexpr\ht\pandoc@box+\dp\pandoc@box\relax}%
  \Gscale@div\@tempb{\linewidth}{\wd\pandoc@box}%
  \ifdim\@tempb\p@<\@tempa\p@\let\@tempa\@tempb\fi% select the smaller of both
  \ifdim\@tempa\p@<\p@\scalebox{\@tempa}{\usebox\pandoc@box}%
  \else\usebox{\pandoc@box}%
  \fi%
}
% Set default figure placement to htbp
\def\fps@figure{htbp}
\makeatother
\setlength{\emergencystretch}{3em} % prevent overfull lines
\providecommand{\tightlist}{%
  \setlength{\itemsep}{0pt}\setlength{\parskip}{0pt}}
\usepackage{bookmark}
\IfFileExists{xurl.sty}{\usepackage{xurl}}{} % add URL line breaks if available
\urlstyle{same}
\hypersetup{
  pdftitle={Introduction to R},
  hidelinks,
  pdfcreator={LaTeX via pandoc}}

\title{Introduction to R}
\author{}
\date{\vspace{-2.5em}}

\begin{document}
\maketitle

\section{Getting Started with R}\label{getting-started-with-r}

Welcome to the R Basics page! We'll walk through installing R, running
simple operations, creating objects, understanding data types, and more.

\subsection{Installing R and RStudio}\label{installing-r-and-rstudio}

\begin{enumerate}
\def\labelenumi{\arabic{enumi}.}
\tightlist
\item
  \textbf{Download R} from the CRAN website:
  \url{https://cran.r-project.org/}
\item
  \textbf{Download RStudio} (a user-friendly IDE):
  \url{https://posit.co/download/rstudio-desktop/}
\item
  Install R first, then RStudio.
\end{enumerate}

\subsection{Interacting with R}\label{interacting-with-r}

To interact with R, we need R-Studio which is very user friendly. By
default, R studio has 4 different panes. The top-left pane shows
R-script where you would like to code and save your work as you go. To
create a new script, go to File \textgreater{} New File \textgreater{} R
Script, or press \texttt{ctrl(command)\ +\ shift\ +\ N}. To run your
commands, press \texttt{ctrl\ +\ Enter} or \texttt{command\ +\ Enter}.
To run all your code at once, press \texttt{Ctrl\ +\ A\ +\ Enter}.

The bottom left window is your console. If you quickly want to check
something or run a code without needing it to be saved on your r-script,
you can directly type into the console. The output/messages/errors are
also shown in the console once you run the code.

Next, the top right window is called ``The Environment''. This window
shows vectors, matrices, dataframes, output etc. in form of objects.

The last window on the bottom right shows plots and graphs. There are
also other tabs such as ``Files'', ``History'', ``Packages and so on.

\subsection{Running Basic Operations}\label{running-basic-operations}

You can use R like a calculator. Try running these in the console, or
any other calculations you would do in the calculator.

Addition

\begin{Shaded}
\begin{Highlighting}[]
\DecValTok{5}\SpecialCharTok{+}\DecValTok{3}
\end{Highlighting}
\end{Shaded}

\begin{verbatim}
## [1] 8
\end{verbatim}

Substraction

\begin{Shaded}
\begin{Highlighting}[]
\DecValTok{5{-}3}
\end{Highlighting}
\end{Shaded}

\begin{verbatim}
## [1] 2
\end{verbatim}

Multiplication

\begin{Shaded}
\begin{Highlighting}[]
\DecValTok{5} \SpecialCharTok{*} \DecValTok{3}
\end{Highlighting}
\end{Shaded}

\begin{verbatim}
## [1] 15
\end{verbatim}

Division

\begin{Shaded}
\begin{Highlighting}[]
\DecValTok{5}\SpecialCharTok{/}\DecValTok{3}
\end{Highlighting}
\end{Shaded}

\begin{verbatim}
## [1] 1.666667
\end{verbatim}

Simplification

\begin{Shaded}
\begin{Highlighting}[]
\DecValTok{5}\SpecialCharTok{+}\DecValTok{3} \SpecialCharTok{{-}} \DecValTok{2}\SpecialCharTok{\^{}}\DecValTok{3}
\end{Highlighting}
\end{Shaded}

\begin{verbatim}
## [1] 0
\end{verbatim}

\subsection{Basic Functions}\label{basic-functions}

R has some basic functions installed already. Some of examples are given
below.

\begin{Shaded}
\begin{Highlighting}[]
\FunctionTok{sum}\NormalTok{(}\FunctionTok{c}\NormalTok{(}\DecValTok{1}\NormalTok{, }\DecValTok{2}\NormalTok{, }\DecValTok{3}\NormalTok{))            }\CommentTok{\# Calculates sum}
\end{Highlighting}
\end{Shaded}

\begin{verbatim}
## [1] 6
\end{verbatim}

\begin{Shaded}
\begin{Highlighting}[]
\FunctionTok{mean}\NormalTok{(}\FunctionTok{c}\NormalTok{(}\DecValTok{4}\NormalTok{, }\DecValTok{5}\NormalTok{, }\DecValTok{6}\NormalTok{))           }\CommentTok{\# Calculates mean}
\end{Highlighting}
\end{Shaded}

\begin{verbatim}
## [1] 5
\end{verbatim}

\begin{Shaded}
\begin{Highlighting}[]
\FunctionTok{seq}\NormalTok{(}\DecValTok{1}\NormalTok{, }\DecValTok{10}\NormalTok{, }\AttributeTok{by =} \DecValTok{2}\NormalTok{)         }\CommentTok{\# Generates a sequence}
\end{Highlighting}
\end{Shaded}

\begin{verbatim}
## [1] 1 3 5 7 9
\end{verbatim}

\begin{Shaded}
\begin{Highlighting}[]
\FunctionTok{sample}\NormalTok{(}\FunctionTok{seq}\NormalTok{(}\DecValTok{1}\SpecialCharTok{:}\DecValTok{6}\NormalTok{), }\AttributeTok{size =} \DecValTok{1}\NormalTok{) }\CommentTok{\# Rolling a die}
\end{Highlighting}
\end{Shaded}

\begin{verbatim}
## [1] 4
\end{verbatim}

\subsection{Installing and Loading
Packages}\label{installing-and-loading-packages}

R is open source; anyone can make their own functions and packages. One
of the most commonly used and versatile package is \texttt{tidyverse}.
You can install and load it as follows. Note that you only need to
install packages once.

\begin{Shaded}
\begin{Highlighting}[]
\FunctionTok{install.packages}\NormalTok{(}\StringTok{"tidyverse"}\NormalTok{) }\CommentTok{\#run this in console}
\FunctionTok{library}\NormalTok{(tidyverse) }
\end{Highlighting}
\end{Shaded}

\subsection{Getting Help}\label{getting-help}

Different help commands allow you to go through the arguments and
details of a particular function. For example, if you want to know how
to use \texttt{mean()} function, you can do it in the following way.

\begin{Shaded}
\begin{Highlighting}[]
\NormalTok{?mean}
\FunctionTok{help}\NormalTok{(mean)}
\end{Highlighting}
\end{Shaded}

\subsection{Creating Objects}\label{creating-objects}

R is an object-based programming. You can assign values to objects using
\texttt{\textless{}-} or \texttt{=} operator. You can assign values,
data, integers, characters et cetera to the object.

\begin{Shaded}
\begin{Highlighting}[]
\NormalTok{a }\OtherTok{\textless{}{-}} \DecValTok{1}
\NormalTok{b }\OtherTok{\textless{}{-}} \DecValTok{2}
\NormalTok{c }\OtherTok{\textless{}{-}} \DecValTok{3}
\end{Highlighting}
\end{Shaded}

We can use \texttt{print()} or just type the object name to view its
value:

\begin{Shaded}
\begin{Highlighting}[]
\NormalTok{a}
\end{Highlighting}
\end{Shaded}

If you are computing something (a simple arithmetic operation or running
a regression itself), you can store those results as well. Stata has
different \texttt{r()} and \texttt{e()} class commands to store results
which can be limiting, but in R, the results can simply be stored into
different objects granting more flexibility.

Note that objects can be overwritten. If I want to assign different
values to a, b and c, I can do so as follows.

\begin{Shaded}
\begin{Highlighting}[]
\NormalTok{a }\OtherTok{\textless{}{-}} \DecValTok{10}
\NormalTok{b }\OtherTok{\textless{}{-}} \DecValTok{30}
\NormalTok{c }\OtherTok{\textless{}{-}} \DecValTok{12}
\end{Highlighting}
\end{Shaded}

If you print these objects, you will get new assigned values instead of
the old ones.

Now, let us suppose a school student asked us to solve quadratic
equations of the form \(ax^2 + bx + c\). The quadratic formula we all
know for this is:

\[
\frac{-b \pm \sqrt{b^2 - 4ac}}{2a}
\] We have defined the values for a, b and c in the previous code
snippet. So, we can simply calculate using the quadratic formula.

\begin{Shaded}
\begin{Highlighting}[]
\NormalTok{(}\SpecialCharTok{{-}}\NormalTok{b }\SpecialCharTok{+} \FunctionTok{sqrt}\NormalTok{(b}\SpecialCharTok{\^{}}\DecValTok{2} \SpecialCharTok{{-}} \DecValTok{4}\SpecialCharTok{*}\NormalTok{a}\SpecialCharTok{*}\NormalTok{c)) }\SpecialCharTok{/}\NormalTok{ (}\DecValTok{2} \SpecialCharTok{*}\NormalTok{ a)}
\end{Highlighting}
\end{Shaded}

\begin{verbatim}
## [1] -0.4753049
\end{verbatim}

\begin{Shaded}
\begin{Highlighting}[]
\NormalTok{(}\SpecialCharTok{{-}}\NormalTok{b }\SpecialCharTok{{-}} \FunctionTok{sqrt}\NormalTok{(b}\SpecialCharTok{\^{}}\DecValTok{2} \SpecialCharTok{{-}} \DecValTok{4}\SpecialCharTok{*}\NormalTok{a}\SpecialCharTok{*}\NormalTok{c)) }\SpecialCharTok{/}\NormalTok{ (}\DecValTok{2} \SpecialCharTok{*}\NormalTok{ a)}
\end{Highlighting}
\end{Shaded}

\begin{verbatim}
## [1] -2.524695
\end{verbatim}

You can store these results in object too.

\begin{Shaded}
\begin{Highlighting}[]
\NormalTok{result }\OtherTok{\textless{}{-}}\NormalTok{ (}\SpecialCharTok{{-}}\NormalTok{b }\SpecialCharTok{+} \FunctionTok{sqrt}\NormalTok{(b}\SpecialCharTok{\^{}}\DecValTok{2} \SpecialCharTok{{-}} \DecValTok{4}\SpecialCharTok{*}\NormalTok{a}\SpecialCharTok{*}\NormalTok{c)) }\SpecialCharTok{/}\NormalTok{ (}\DecValTok{2} \SpecialCharTok{*}\NormalTok{ a)}
\FunctionTok{print}\NormalTok{(result)}
\end{Highlighting}
\end{Shaded}

\begin{verbatim}
## [1] -0.4753049
\end{verbatim}

Sometimes we might need to use different values of a,b and c to
calculate the quadratic equation. For such repetitive tasks, we can make
our own functions and run loops to compute at once. We'll talk about it
later on in the camp.

\subsection{Classes of Objects}\label{classes-of-objects}

In STATA or excel, some columns are numbers and some can be characters/
strings. They could also be boolean or logical (\texttt{TRUE} or
\texttt{FALSE}). R has different terminologies and handles the data or
objects based on the type it is defined as. We call it \texttt{class()}
in R.

Some common types of classes are as follows:

\begin{Shaded}
\begin{Highlighting}[]
\NormalTok{num }\OtherTok{\textless{}{-}} \FloatTok{3.14}       \CommentTok{\# numeric}
\NormalTok{int }\OtherTok{\textless{}{-}} \DecValTok{42}\DataTypeTok{L}        \CommentTok{\# integer}
\NormalTok{char }\OtherTok{\textless{}{-}} \StringTok{"Hello"}   \CommentTok{\# character}
\NormalTok{logic }\OtherTok{\textless{}{-}} \ConstantTok{TRUE}     \CommentTok{\# logical}
\NormalTok{gender }\OtherTok{\textless{}{-}} \FunctionTok{factor}\NormalTok{(}\StringTok{"male"}\NormalTok{, }\StringTok{"female"}\NormalTok{) }\CommentTok{\# factor}
\end{Highlighting}
\end{Shaded}

You can check the class of objects using \texttt{class} or
\texttt{typeof} command.

\begin{Shaded}
\begin{Highlighting}[]
\FunctionTok{class}\NormalTok{(num)}
\FunctionTok{class}\NormalTok{(char)}
\end{Highlighting}
\end{Shaded}

\subsection{Vectors}\label{vectors}

Vectors are a basic data structure in R. A single number is technically
a vector of length 1. The function \texttt{length} tells us how many
entries are there in a vector.

\subsubsection{Numeric}\label{numeric}

Lets create a vector of our own. Suppose we have 4 graduate students in
a classroom and we want to store their age in years.

\begin{Shaded}
\begin{Highlighting}[]
\NormalTok{age }\OtherTok{\textless{}{-}} \FunctionTok{c}\NormalTok{(}\DecValTok{20}\NormalTok{, }\DecValTok{25}\NormalTok{, }\DecValTok{30}\NormalTok{, }\DecValTok{28}\NormalTok{)}
\FunctionTok{length}\NormalTok{(age)}
\end{Highlighting}
\end{Shaded}

\begin{verbatim}
## [1] 4
\end{verbatim}

We see that the vector has 4 entries. But what is the class of this
vector? Lets check.

\begin{Shaded}
\begin{Highlighting}[]
\FunctionTok{class}\NormalTok{(age)}
\end{Highlighting}
\end{Shaded}

\begin{verbatim}
## [1] "numeric"
\end{verbatim}

Similarly, lets create a vector of income of students.

\begin{Shaded}
\begin{Highlighting}[]
\NormalTok{income }\OtherTok{\textless{}{-}} \FunctionTok{c}\NormalTok{(}\DecValTok{5}\NormalTok{,}\DecValTok{10}\NormalTok{,}\DecValTok{20}\NormalTok{,}\DecValTok{30}\NormalTok{)}
\FunctionTok{class}\NormalTok{(income)}
\end{Highlighting}
\end{Shaded}

\begin{verbatim}
## [1] "numeric"
\end{verbatim}

\subsubsection{Factors (Categorical
Variables)}\label{factors-categorical-variables}

Factors denote a class for storing categorical data in R. Suppose, now
we also want to store the information on gender of 4 students in the
classroom.

\begin{Shaded}
\begin{Highlighting}[]
\NormalTok{gender }\OtherTok{\textless{}{-}} \FunctionTok{factor}\NormalTok{(}\FunctionTok{c}\NormalTok{(}\StringTok{"male"}\NormalTok{, }\StringTok{"female"}\NormalTok{, }\StringTok{"female"}\NormalTok{, }\StringTok{"male"}\NormalTok{))}
\end{Highlighting}
\end{Shaded}

\subsubsection{Logical}\label{logical}

Now, we want to know whether these students attended the code camp at
CSU. Lets store that information as well as a vector.

\begin{Shaded}
\begin{Highlighting}[]
\NormalTok{attend\_camp }\OtherTok{\textless{}{-}} \FunctionTok{c}\NormalTok{(}\ConstantTok{TRUE}\NormalTok{, }\ConstantTok{FALSE}\NormalTok{, }\ConstantTok{FALSE}\NormalTok{, }\ConstantTok{TRUE}\NormalTok{)}
\end{Highlighting}
\end{Shaded}

Note that, normally in statistical analysis, we store this information
as 1 or 0. But I am using logical vector here just for the illustration
purposes.

\subsection{Data Frames}\label{data-frames}

Until now, we just worked with numbers and vectors. In reality, as
economists, we would have to work with large dataset. Like excel
spreadsheet, R has an object class \texttt{dataframe} that stores the
data in a spreadsheet-like format.

Now, lets try combing the information on 4 students in a single data
frame.

\begin{Shaded}
\begin{Highlighting}[]
\NormalTok{dat\_grads }\OtherTok{\textless{}{-}} \FunctionTok{data.frame}\NormalTok{(age, gender, income, attend\_camp, }\AttributeTok{id =} \FunctionTok{seq}\NormalTok{(}\DecValTok{1}\SpecialCharTok{:}\DecValTok{4}\NormalTok{))}
\FunctionTok{head}\NormalTok{(dat\_grads)}
\end{Highlighting}
\end{Shaded}

\begin{verbatim}
##   age gender income attend_camp id
## 1  20   male      5        TRUE  1
## 2  25 female     10       FALSE  2
## 3  30 female     20       FALSE  3
## 4  28   male     30        TRUE  4
\end{verbatim}

\subsection{The accessor}\label{the-accessor}

The accessor \texttt{\$} is used to access particular columns within a
dataframe in R.

\begin{Shaded}
\begin{Highlighting}[]
\CommentTok{\# What is the frequency of male and female?}
\FunctionTok{table}\NormalTok{(dat\_grads}\SpecialCharTok{$}\NormalTok{gender)}
\end{Highlighting}
\end{Shaded}

\begin{verbatim}
## 
## female   male 
##      2      2
\end{verbatim}

\subsection{Matrices}\label{matrices}

Matrix is very important in econometric analysis. All the econometrics
commands in STATA, R or any other statistical packages are based on
matrix. If you have ever heard or used GAUSS, you will need to do
regressions and stuff using the matrix operations. Thanks to developers
these days who make pre-built packages where we could just input our
data frame and variables and get regression results.

Lets start off with a simple \(2\times2\) matrix.

\begin{Shaded}
\begin{Highlighting}[]
\NormalTok{mat1 }\OtherTok{\textless{}{-}} \FunctionTok{matrix}\NormalTok{(}\FunctionTok{c}\NormalTok{(}\DecValTok{1}\NormalTok{,}\DecValTok{2}\NormalTok{,}\DecValTok{3}\NormalTok{,}\DecValTok{4}\NormalTok{), }\AttributeTok{nrow =} \DecValTok{2}\NormalTok{, }\AttributeTok{ncol =} \DecValTok{2}\NormalTok{, }\AttributeTok{byrow =} \ConstantTok{TRUE}\NormalTok{)}
\FunctionTok{print}\NormalTok{(mat1)}
\end{Highlighting}
\end{Shaded}

\begin{verbatim}
##      [,1] [,2]
## [1,]    1    2
## [2,]    3    4
\end{verbatim}

Similarly, lets create a \(2\times3\) matrix.

\begin{Shaded}
\begin{Highlighting}[]
\NormalTok{mat2 }\OtherTok{\textless{}{-}} \FunctionTok{matrix}\NormalTok{(}\FunctionTok{c}\NormalTok{(}\DecValTok{1}\NormalTok{,}\DecValTok{2}\NormalTok{,}\DecValTok{3}\NormalTok{,}\DecValTok{4}\NormalTok{,}\DecValTok{5}\NormalTok{,}\DecValTok{6}\NormalTok{), }\AttributeTok{nrow =} \DecValTok{2}\NormalTok{, }\AttributeTok{ncol =} \DecValTok{6}\NormalTok{)}
\end{Highlighting}
\end{Shaded}

Lets do some matrix operations which we commonly use including addition,
subtraction, multiplication, inverse, transpose et. cetera. Note that
the dimensions need to match following matrix rules, otherwise R will
return error.

\begin{Shaded}
\begin{Highlighting}[]
\NormalTok{mat1 }\SpecialCharTok{+}\NormalTok{ mat1 }\CommentTok{\# Addition}
\end{Highlighting}
\end{Shaded}

\begin{verbatim}
##      [,1] [,2]
## [1,]    2    4
## [2,]    6    8
\end{verbatim}

\begin{Shaded}
\begin{Highlighting}[]
\NormalTok{mat1 }\SpecialCharTok{\%*\%}\NormalTok{ mat2 }\CommentTok{\# Multiplication}
\end{Highlighting}
\end{Shaded}

\begin{verbatim}
##      [,1] [,2] [,3] [,4] [,5] [,6]
## [1,]    5   11   17    5   11   17
## [2,]   11   25   39   11   25   39
\end{verbatim}

\begin{Shaded}
\begin{Highlighting}[]
\FunctionTok{t}\NormalTok{(mat2)     }\CommentTok{\#Transpose}
\end{Highlighting}
\end{Shaded}

\begin{verbatim}
##      [,1] [,2]
## [1,]    1    2
## [2,]    3    4
## [3,]    5    6
## [4,]    1    2
## [5,]    3    4
## [6,]    5    6
\end{verbatim}

\begin{Shaded}
\begin{Highlighting}[]
\FunctionTok{solve}\NormalTok{(mat1) }\CommentTok{\# Inverse}
\end{Highlighting}
\end{Shaded}

\begin{verbatim}
##      [,1] [,2]
## [1,] -2.0  1.0
## [2,]  1.5 -0.5
\end{verbatim}

\subsection{Comment your work!}\label{comment-your-work}

Commenting your code is important for two main reasons. First, as a
scientist, you should strive to make your work reproducible. Annotating
your code makes it easier to understand for other people to interpret
your process. Second, you will inevitably write a lot of code, get busy,
move on to other things, and reopen your code months later with little
to no memory of what you did. Spending the extra 5 minutes to annotate
your work will save you a headache down the line.

To comment your code, simply put a hashtag character, \#, in front of
your annotations. R treats \# in a special way, and it will not run
anything that follows \# on a line.

\begin{Shaded}
\begin{Highlighting}[]
\CommentTok{\# I create an object first}
\NormalTok{my\_first\_object }\OtherTok{\textless{}{-}} \DecValTok{3}


\CommentTok{\# I double my object to create a second object}
\NormalTok{my\_second\_object }\OtherTok{\textless{}{-}}\NormalTok{ my\_first\_object }\SpecialCharTok{*} \DecValTok{2}
\end{Highlighting}
\end{Shaded}

\subsection{Indexing}\label{indexing}

Before we dive into R functions, lets learn about indexing. Indexing
refers to the process of selecting specific elements or subsets of data
from vectors, matrices, dataframes or lists using their positions
(indices) or other criteria. In R, we use square bracket \texttt{{[}{]}}
for indexing.

\subsubsection{Integer Indexing}\label{integer-indexing}

Lets say we want to see what specific element of a vector looks like.

\begin{Shaded}
\begin{Highlighting}[]
\CommentTok{\# What is the 2nd element in the age vector?}
\NormalTok{age[}\DecValTok{2}\NormalTok{]}
\end{Highlighting}
\end{Shaded}

\begin{verbatim}
## [1] 25
\end{verbatim}

\subsubsection{Negative Integer
Indexing}\label{negative-integer-indexing}

If you want all the elements except 2nd element, we can use negative
indexing.

\begin{Shaded}
\begin{Highlighting}[]
\NormalTok{age[}\SpecialCharTok{{-}}\DecValTok{2}\NormalTok{]}
\end{Highlighting}
\end{Shaded}

\begin{verbatim}
## [1] 20 30 28
\end{verbatim}

If we want to access multiple elements of a vector at once, we can use
concatenate function \texttt{c()}.

\begin{Shaded}
\begin{Highlighting}[]
\CommentTok{\# What is the 3rd and 4th element in the age vector?}
\NormalTok{age[}\FunctionTok{c}\NormalTok{(}\DecValTok{3}\NormalTok{,}\DecValTok{4}\NormalTok{)]}
\end{Highlighting}
\end{Shaded}

\begin{verbatim}
## [1] 30 28
\end{verbatim}

Indexing works not only in vectors, but also in matrices, dataframes and
other classes of objects. Lets say we want to access first column of our
dataset.

\begin{Shaded}
\begin{Highlighting}[]
\CommentTok{\# Print the first column of our graduate student dataset}
\NormalTok{dat\_grads[, }\DecValTok{1}\NormalTok{]}
\end{Highlighting}
\end{Shaded}

\begin{verbatim}
## [1] 20 25 30 28
\end{verbatim}

Now, instead of column, suppose we want to print data of a particular
respondent/id (i.e., print rows instead of columns).

\begin{Shaded}
\begin{Highlighting}[]
\CommentTok{\# Print the data of 3rd and 4th student in the dataset}
\NormalTok{dat\_grads[}\FunctionTok{c}\NormalTok{(}\DecValTok{3}\NormalTok{,}\DecValTok{4}\NormalTok{), ]}
\end{Highlighting}
\end{Shaded}

\begin{verbatim}
##   age gender income attend_camp id
## 3  30 female     20       FALSE  3
## 4  28   male     30        TRUE  4
\end{verbatim}

\subsubsection{Logical Indexing}\label{logical-indexing}

Suppose you want to filter your data using a condition. For example, we
want a dataset of the students who have income more than \$10.

\begin{Shaded}
\begin{Highlighting}[]
\CommentTok{\# Print the data of students with income more than $10}
\NormalTok{dat\_grads[dat\_grads}\SpecialCharTok{$}\NormalTok{income }\SpecialCharTok{\textgreater{}} \DecValTok{10}\NormalTok{, ]}
\end{Highlighting}
\end{Shaded}

\begin{verbatim}
##   age gender income attend_camp id
## 3  30 female     20       FALSE  3
## 4  28   male     30        TRUE  4
\end{verbatim}

\subsubsection{Name Indexing}\label{name-indexing}

We can also use names or characters to select columns or filter
elements.

\begin{Shaded}
\begin{Highlighting}[]
\CommentTok{\# Print data of students with id = 3 and 4}
\NormalTok{dat\_grads[dat\_grads}\SpecialCharTok{$}\NormalTok{id }\SpecialCharTok{\%in\%} \FunctionTok{c}\NormalTok{(}\DecValTok{3}\NormalTok{,}\DecValTok{4}\NormalTok{), ]}
\end{Highlighting}
\end{Shaded}

\begin{verbatim}
##   age gender income attend_camp id
## 3  30 female     20       FALSE  3
## 4  28   male     30        TRUE  4
\end{verbatim}

\subsection{Basic Summary Statistics using R
functions}\label{basic-summary-statistics-using-r-functions}

Now we can calculate some descriptive statistics such as mean, standard
deviation and so on. We will learn about summary statistics for
dataframes in the next section. For now, let us just focus on vector.

\begin{Shaded}
\begin{Highlighting}[]
\CommentTok{\# Creating a random vector. Also include the elements of vector age in this new vector}
\NormalTok{my\_third\_object }\OtherTok{\textless{}{-}} \FunctionTok{c}\NormalTok{(}\DecValTok{23}\NormalTok{, }\DecValTok{24}\NormalTok{, age, }\DecValTok{26}\NormalTok{, }\DecValTok{30}\NormalTok{, }\DecValTok{34}\NormalTok{, }\DecValTok{36}\NormalTok{, }\DecValTok{40}\NormalTok{)}

\CommentTok{\# What is the length of the vector?}
\FunctionTok{length}\NormalTok{(my\_third\_object)}
\end{Highlighting}
\end{Shaded}

\begin{verbatim}
## [1] 11
\end{verbatim}

Now, lets calculate mean, median, standard deviation, minimum and
maximum of our vector.

\begin{Shaded}
\begin{Highlighting}[]
\CommentTok{\# Minimum value}
\FunctionTok{min}\NormalTok{(my\_third\_object)}
\end{Highlighting}
\end{Shaded}

\begin{verbatim}
## [1] 20
\end{verbatim}

\begin{Shaded}
\begin{Highlighting}[]
\CommentTok{\# Maximum value}
\FunctionTok{max}\NormalTok{(my\_third\_object)}
\end{Highlighting}
\end{Shaded}

\begin{verbatim}
## [1] 40
\end{verbatim}

\begin{Shaded}
\begin{Highlighting}[]
\CommentTok{\# Mean}
\FunctionTok{mean}\NormalTok{(my\_third\_object)}
\end{Highlighting}
\end{Shaded}

\begin{verbatim}
## [1] 28.72727
\end{verbatim}

\begin{Shaded}
\begin{Highlighting}[]
\CommentTok{\# Median}
\FunctionTok{median}\NormalTok{(my\_third\_object)}
\end{Highlighting}
\end{Shaded}

\begin{verbatim}
## [1] 28
\end{verbatim}

\begin{Shaded}
\begin{Highlighting}[]
\CommentTok{\# Standard deviation}
\FunctionTok{sd}\NormalTok{(my\_third\_object)}
\end{Highlighting}
\end{Shaded}

\begin{verbatim}
## [1] 6.034748
\end{verbatim}

You can also use \texttt{summary()} function.

Now, lets calculate mean manually and check if it matches with the R
in-built commands. You can do the same for other statistics on your own.

\begin{Shaded}
\begin{Highlighting}[]
\CommentTok{\# Calculating mean manually}
\NormalTok{my\_mean }\OtherTok{\textless{}{-}} \FunctionTok{sum}\NormalTok{(my\_third\_object) }\SpecialCharTok{/} \FunctionTok{length}\NormalTok{(my\_third\_object)}

\CommentTok{\# Check if this result matches the one calculated by R\textquotesingle{}s in{-}built function}
\FunctionTok{mean}\NormalTok{(my\_third\_object) }\SpecialCharTok{==}\NormalTok{ my\_mean}
\end{Highlighting}
\end{Shaded}

\begin{verbatim}
## [1] TRUE
\end{verbatim}

\subsection{Some Other functions}\label{some-other-functions}

Lets talk about some in-built functions that we frequently use. We will
also go through some examples of nesting. The functions in R can be
nested within the other function.

Lets say we want to know the number of unique observations in our
vector. We first use \texttt{unique()} function to print unique values
in the vector and next it within \texttt{length()} function to give
number of unique values.

\begin{Shaded}
\begin{Highlighting}[]
\FunctionTok{length}\NormalTok{(}\FunctionTok{unique}\NormalTok{(my\_third\_object))}
\end{Highlighting}
\end{Shaded}

\begin{verbatim}
## [1] 10
\end{verbatim}

Similarly, lets calculate the mean again but now, we want to round it
off to 2 digits.

\begin{Shaded}
\begin{Highlighting}[]
\FunctionTok{round}\NormalTok{(}\FunctionTok{mean}\NormalTok{(my\_third\_object), }\AttributeTok{digits =} \DecValTok{2}\NormalTok{)}
\end{Highlighting}
\end{Shaded}

\begin{verbatim}
## [1] 28.73
\end{verbatim}

Now, lets sample some elements of the vector. Each time you sample, you
get different results.

\begin{Shaded}
\begin{Highlighting}[]
\CommentTok{\# Sample 5 elements}
\FunctionTok{sample}\NormalTok{(my\_third\_object, }\AttributeTok{size =} \DecValTok{5}\NormalTok{, }\AttributeTok{replace =} \ConstantTok{TRUE}\NormalTok{)}
\end{Highlighting}
\end{Shaded}

\begin{verbatim}
## [1] 24 26 26 24 28
\end{verbatim}

\begin{Shaded}
\begin{Highlighting}[]
\CommentTok{\# Sample 5 elements again}
\FunctionTok{sample}\NormalTok{(my\_third\_object, }\AttributeTok{size =} \DecValTok{5}\NormalTok{, }\AttributeTok{replace =} \ConstantTok{TRUE}\NormalTok{)}
\end{Highlighting}
\end{Shaded}

\begin{verbatim}
## [1] 36 40 36 30 28
\end{verbatim}

Anytime your code involves random processes (such as sampling), include
the \texttt{set.seed()} function at the beginning of your script to make
your results reproducible. That way, the next time you run your code,
the same samples will be generated from when you first wrote the script.

\begin{Shaded}
\begin{Highlighting}[]
\FunctionTok{set.seed}\NormalTok{(}\DecValTok{12345}\NormalTok{)}
\end{Highlighting}
\end{Shaded}

\subsubsection{Writing your own
functions}\label{writing-your-own-functions}

You will often find yourself in situations where you have to calculate
same thing repetitively for different data or variables. In such
situations, writing your own functions as well as running loops will
allow you to get your desired results efficiently.

\begin{Shaded}
\begin{Highlighting}[]
\CommentTok{\# Write a function that calculates circumference of a circle given a radius}
\NormalTok{circumference }\OtherTok{\textless{}{-}} \ControlFlowTok{function}\NormalTok{(r) \{}
\NormalTok{  result }\OtherTok{\textless{}{-}} \DecValTok{2} \SpecialCharTok{*}\NormalTok{ pi }\SpecialCharTok{*}\NormalTok{ r}
  \FunctionTok{print}\NormalTok{(result)}
\NormalTok{\}}
\end{Highlighting}
\end{Shaded}

\begin{Shaded}
\begin{Highlighting}[]
\FunctionTok{circumference}\NormalTok{(}\AttributeTok{r =} \DecValTok{2}\NormalTok{)}
\end{Highlighting}
\end{Shaded}

\begin{verbatim}
## [1] 12.56637
\end{verbatim}

\begin{Shaded}
\begin{Highlighting}[]
\FunctionTok{circumference}\NormalTok{(}\AttributeTok{r =} \DecValTok{10}\NormalTok{)}
\end{Highlighting}
\end{Shaded}

\begin{verbatim}
## [1] 62.83185
\end{verbatim}

Lets try another example. Suppose you roll two die and you want to find
out the sum of the outcome.

\begin{Shaded}
\begin{Highlighting}[]
\NormalTok{roll }\OtherTok{\textless{}{-}} \ControlFlowTok{function}\NormalTok{() \{}
\NormalTok{  die }\OtherTok{\textless{}{-}} \FunctionTok{sample}\NormalTok{(}\FunctionTok{seq}\NormalTok{(}\DecValTok{1}\SpecialCharTok{:}\DecValTok{6}\NormalTok{), }\AttributeTok{replace =} \ConstantTok{TRUE}\NormalTok{, }\AttributeTok{size =} \DecValTok{2}\NormalTok{)}
  \FunctionTok{return}\NormalTok{(}\FunctionTok{sum}\NormalTok{(die))}
\NormalTok{\}}
\end{Highlighting}
\end{Shaded}

Now try rolling the die and notice that you get different outcomes
everytime.

\begin{Shaded}
\begin{Highlighting}[]
\FunctionTok{roll}\NormalTok{()}
\end{Highlighting}
\end{Shaded}

\begin{verbatim}
## [1] 9
\end{verbatim}

\begin{Shaded}
\begin{Highlighting}[]
\FunctionTok{roll}\NormalTok{()}
\end{Highlighting}
\end{Shaded}

\begin{verbatim}
## [1] 6
\end{verbatim}

\begin{Shaded}
\begin{Highlighting}[]
\FunctionTok{roll}\NormalTok{()}
\end{Highlighting}
\end{Shaded}

\begin{verbatim}
## [1] 7
\end{verbatim}

\subsubsection{For loops}\label{for-loops}

For loops are used to iterate a set of operations over a collection of
objects such as vector, data frame, matrix or lists.

\begin{Shaded}
\begin{Highlighting}[]
\ControlFlowTok{for}\NormalTok{ (variable }\ControlFlowTok{in}\NormalTok{ sequence) \{}
\NormalTok{    expression}
\NormalTok{\}}
\end{Highlighting}
\end{Shaded}

Suppose you have a vector of radius and you need to calculate
circumference for each one of those. Instead of computing circumference
separately for each radius, we can use for loop to conduct such
repetitive task.

\begin{Shaded}
\begin{Highlighting}[]
\CommentTok{\# Create a vector of radius}
\NormalTok{radius }\OtherTok{\textless{}{-}} \FunctionTok{c}\NormalTok{(}\DecValTok{1}\NormalTok{,}\DecValTok{3}\NormalTok{,}\DecValTok{4}\NormalTok{,}\DecValTok{6}\NormalTok{,}\DecValTok{8}\NormalTok{,}\DecValTok{9}\NormalTok{,}\DecValTok{10}\NormalTok{)}

\CommentTok{\# Loop accross all radius values to compute circumference}
\ControlFlowTok{for}\NormalTok{ (i }\ControlFlowTok{in}\NormalTok{ radius) \{}
  \FunctionTok{circumference}\NormalTok{(}\AttributeTok{r =}\NormalTok{ i)}
\NormalTok{\}}
\end{Highlighting}
\end{Shaded}

\begin{verbatim}
## [1] 6.283185
## [1] 18.84956
## [1] 25.13274
## [1] 37.69911
## [1] 50.26548
## [1] 56.54867
## [1] 62.83185
\end{verbatim}

\begin{center}\rule{0.5\linewidth}{0.5pt}\end{center}

\end{document}
